% !TEX root = ../main.tex
\chapter{Literature Review and Theoretical Framework}
\label{ch:literature}

My research problem in Chapter~\ref{ch:introduction} is that transfer-graph reputation scores conflate payment flows with trust. My objectives O1--O5 and my research questions RQ1--RQ5 emerge from that problem. I compare EndorseRank and Adaptive Weighted PageRank (AWP) on that problem by running their graphs through PageRank and checking the rankings using rank correlations. The framework comprises hyperlink--citation propagation, domain alignment and computational complexity. I take up the remaining gap in Chapters~\ref{ch:methodology} and~\ref{ch:results}.

\dissertationSubheading[sec:blockchain-defi-foundations]{Blockchain and DeFi foundations}

A public blockchain records transactions in hashed blocks that later blocks secure \cend{5}. After a block has been approved by a majority of the participating nodes, changing it becomes costly, enabling parties to settle on the ledger without requiring each party to have a known legal identity (\emph{trust-minimized} settlement).

Bitcoin used this ledger for a peer-to-peer electronic cash network \cend{5}. Ethereum retains the same append-only ledger, but allows an account to store a program as well as a balance, so a transaction can invoke a function, modifying the contract's state \cend{7}; those functions are executed on the Ethereum Virtual Machine (EVM) \cend{36}. The programs resident on the ledger are termed smart contracts, and any honest node replaying a function invocation will produce the same output \cend{6}. This allows collateralization ratios, liquidation conditions and fee schedules to be enforced by the contract instead of a trusted third party \cend{1}.

These contracts are known collectively as Decentralized Finance (DeFi). They provide the functionality of exchanges, lending platforms, derivatives markets and asset management services; Werner et al.\cend{8} define the protocol classes, and Sch\"{a}r\cend{1} outlines the layering of protocols. The ability of one protocol to interact with another allows new financial instruments to be composed out of existing ones \cend{37}, but also means that failure can propagate across multiple protocols. Examples of successful attacks include flash-loan exploits \cend{38} and price-oracle manipulation \cend{39}. Transaction histories in these chains are publicly available and machine-readable, but they reflect actions, not legal identities \cend{2}.

\dissertationSubsubheading[sub:pseudonymity-credit]{Pseudonymity and the credit problem}

In conventional finance, credit scoring associates a legal entity with a credit history. On a public blockchain, the primary unit of analysis is an address (a wallet), which may represent a legal entity, a trust, an algorithmic agent, or a smart contract. A single entity can own multiple addresses \cend{11}, and statistical clustering techniques can identify ownership patterns only partially \cend{40}. As a result, in on-chain lending, borrowers must post more value in collateral than the amount they borrow, and liquidators are permitted to take ownership of the excess collateral if the borrower falls below a certain threshold \cend{41}. This approach is secure but capital-inefficient. On-chain credit scores attempt to quantify the likelihood that an actor will repay a loan, allowing a protocol designer, or a researcher attempting to study credit risk, to evaluate actors based on on-chain behavior without requiring access to off-chain identity information. Recent research has proposed credit scoring mechanisms using prospect theory \cend{13}, applying machine learning to on-chain lending data \cend{14} and modeling credit risk as a dynamic process propagating through account networks \cend{12}.

\dissertationSubheading[sec:ethereum-l2]{Ethereum, Layer-2 rollups, and Arbitrum One}

Most DeFi applications have deployed on the Ethereum mainnet, as it has historically been the primary venue for DeFi innovation and liquidity, but growing adoption has increased gas costs and constrained throughput \cend{42}. Layer-2 (L2) protocols attempt to maintain Ethereum's security guarantees while reducing the cost of transactions, either by executing transactions off-chain or by batching transactions together and submitting a succinct summary of the transaction execution to the mainnet Ethereum blockchain \cend{43}. Thibault et al.\cend{44} present a taxonomy of L2 scaling approaches. Optimistic L2s, like Arbitrum, assume that submitted batches of transactions are valid, unless there is evidence presented in a dispute period \cend{45}.

Arbitrum One is an example of such an optimistic L2 solution that supports execution of EVM-based smart contracts. It has become a popular platform for DeFi applications, including GMX, a perpetual swaps exchange platform \cend{46}, whose smart contracts have been the most-authorized recipients in our observation period. While perpetuals are a distinct category of derivatives from the fixed-maturity futures studied by Ackerer et al.\cend{47}, they share similarities in the mechanics of funding-rate pricing, which Ackerer et al.\cend{47} model. We use Arbitrum One for this research because both its ERC-20 \texttt{Approval} and \texttt{Transfer} events are included in the publicly available BigQuery dataset, and its fees are low enough that the amount of activity in a six-month period is dense. Note that we do not claim our results will generalize to other L2s or to Ethereum mainnet without further validation, but we believe the framework we propose (approvals and transfers, proxy contracts, and rank-based centrality correlation) applies on any chain that emits the same logs.

\dissertationSubheading[sec:erc20-allowance-lifecycle]{ERC-20 tokens and the allowance lifecycle}

The ERC-20 standard provides a common interface for fungible tokens: \texttt{transfer}, \texttt{approve}, \texttt{transferFrom}, \texttt{allowance}, and associated events \cend{18}. When the owner calls \texttt{approve(spender, amount)}, the token contract stores that the spender can transfer up to \texttt{amount} from the owner's balance. Later, the spender can call \texttt{transferFrom} to execute a token transfer. EIP-2612 \texttt{permit} allows allowance updates using an off-chain signature that can be settled in a single transaction \cend{19}, offering a better user experience without altering the underlying semantics of authorization.

Figure~\ref{fig:allowance-sequence} shows the ERC-20 allowance lifecycle. An allowance represents an authorization, not an actual transfer. Token owners can grant allowances to routers, vaults, or trading contracts they trust to operate within defined limits. Revocation (approving zero) and updating (setting a new non-zero allowance) change the active state of the allowance. EndorseRank takes this into account by considering the most recent allowance value for each owner--spender pair as the edge weight: prior allowances are replaced rather than gradually decaying over time.

\begin{figure}[htbp]
\centering
\begin{tikzpicture}[
  font=\small,
  actor/.style={draw, rounded corners, minimum width=2.4cm, minimum height=0.75cm, align=center},
  msg/.style={-{Stealth[length=2.5mm]}, thick},
  note/.style={font=\scriptsize\itshape, text=black!70, align=left}
]
% Actors
\node[actor] (owner) at (0,0) {Owner};
\node[actor] (token) at (5.5,0) {ERC-20 token};
\node[actor] (spender) at (11,0) {Spender};

% Lifelines (vertical dashed lines)
\draw[dashed, gray] (owner.south) -- ++(0,-5.2);
\draw[dashed, gray] (token.south) -- ++(0,-5.2);
\draw[dashed, gray] (spender.south) -- ++(0,-5.2);

% Message 1: approve (owner -> token), horizontal span 5.5cm
\draw[msg] (0,-1.2) -- node[above,font=\scriptsize] {\texttt{approve(spender,$a$)}} (5.5,-1.2);
\node[note] at (2.75,-1.7) {emit \texttt{Approval}; store allowance};

% Message 2: transferFrom (spender -> token), horizontal span 5.5cm
\draw[msg] (11,-2.8) -- node[above,font=\scriptsize] {\texttt{transferFrom(owner,to,$v\leq a$)}} (5.5,-2.8);
\node[note] at (8.25,-3.3) {move tokens; reduce allowance};

% Message 3: revoke (owner -> token), dashed
\draw[msg, dashed] (0,-4.4) -- node[above,font=\scriptsize] {\texttt{approve(spender,$0$)} (revoke)} (5.5,-4.4);
\end{tikzpicture}
\caption{ERC-20 allowance lifecycle: approve (or permit), optional \texttt{transferFrom}, and revoke. EndorseRank edges reflect the latest live allowance snapshot per owner--spender pair (historical approvals superseded, not time-decayed).}
\label{fig:allowance-sequence}
\end{figure}

From a data-availability standpoint, allowances are sparse compared to transfers: numerous transfers can occur under one approval, while many wallets have never granted approval to external accounts. This sparsity offers both a computational benefit and a semantic distinction: EndorseRank quantifies authorization to spend, while a transfer graph reflects payment flows.

\dissertationSubheading[sec:graph-reputation]{Graph-based on-chain reputation}

PageRank and its variants rank nodes based on influence propagated along directed edges \cend{15}. Two recent studies have applied this approach to transaction graphs on blockchains, which should be distinguished. Do and Do\cend{4} use PageRank and HodgeRank on Ethereum transaction graphs to measure social credit; their main contribution is defining the connectivity between accounts in the context of credit scoring. Adaptive Weighted PageRank (AWP), proposed by Do, Do, and Nguyen\cend{3}, serves as the transfer-graph baseline in this thesis: it assigns each transfer edge a weight based on value and a logistic decay function of edge age, then evaluates the ranking scores against in-degree and in-value (total count and sum of incoming transfers) as proxies for financial significance. The specific form of the logistic decay function and the evaluation methodology with in-degree/in-value metrics used in Chapter~\ref{ch:methodology} are taken from the AWP paper\cend{3}, not from Do and Do\cend{4}. In addition, Lin et al.\cend{12} argue that transaction connectivity may fail to capture structure relevant to risk unless risk is propagated through the network, which justifies more sophisticated graph semantics than vanilla PageRank. Similarly, Wu et al.\cend{24} make the same point for phishing detection, noting that network embeddings of transaction graphs do better than per-account features. In surveys of blockchain-based reputation systems \cend{25}, as well as of smart-contract reputation designs \cend{26}, it is documented that most existing deployments still depend on explicit ratings or transfer volume as the trust signal; a survey of knowledge discovery in cryptocurrency transaction data \cend{48} also lists degree- and value-based centralities as the primary descriptors of wallets. Packin and Lev-Aretz\cend{2} emphasize that systems must be auditable even though their users are pseudonymous. Because EndorseRank relies solely on the propagation of graph signals, it can be audited in the manner prescribed.

Graph-based methods are appealing because they (i) make use of only public information, (ii) yield continuous scores appropriate for ranking, and (iii) allow complexity analysis in terms of number of edges. On the other hand, graph-based methods are susceptible to Sybil attacks \cend{11}, wash trading, and the mismatch between graph centrality and actual default risk. Note that PageRank scores are not credit scores in the regulatory sense; rather, the aim here is simply to compare allowance-based and transfer-based rankings rigorously using the same set of proxies.

\dissertationSubheading[sec:pagerank-formalism]{PageRank formalism}

Let $G=(V,E)$ be a directed graph, with $|V|=n$ nodes. Let $w_{ij}>0$ be the weight of the edge from $i$ to $j$ (representing influence from $i$ to $j$ in the endorsement or transfer sense described above). Let $W$ be the corresponding weighted adjacency matrix and let $d_i^{\mathrm{out}}=\sum_j w_{ij}$. Define the row-stochastic transition matrix $P$ as
\begin{equation}
P_{ij} =
\begin{cases}
w_{ij}/d_i^{\mathrm{out}} & \text{if } d_i^{\mathrm{out}}>0,\\
1/n & \text{if } d_i^{\mathrm{out}}=0
\end{cases}
\label{eq:transition}
\end{equation}
(dangling nodes distribute their mass uniformly). For a damping parameter $d\in(0,1)$, define the Google matrix
\begin{equation}
G_{\mathrm{PR}} = d\, P^{\top} + \frac{1-d}{n}\,\mathbf{1}\mathbf{1}^{\top},
\label{eq:google-matrix}
\end{equation}
and let $\boldsymbol{\pi}$ be the PageRank vector, satisfying $\boldsymbol{\pi}=G_{\mathrm{PR}}\boldsymbol{\pi}$ with $\boldsymbol{\pi}\geq 0$ and $\sum_i \pi_i=1$ \cend{15,23}. In element-wise form, this is
\begin{equation}
\pi_j = d\sum_{i:(i\to j)\in E}\pi_i \frac{w_{ij}}{d_i^{\mathrm{out}}} + \frac{1-d}{n}.
\label{eq:pagerank-node}
\end{equation}

\dissertationSubsubheading[sub:random-surfer]{Random-surfer interpretation}

Equation~\eqref{eq:pagerank-node} has a probabilistic interpretation: at each step, a surfer follows an outgoing edge with probability $d$ and teleports uniformly with probability $1-d$ (Figure~\ref{fig:random-surfer}). The stationary distribution then provides a ranking according to long-run visit frequency. Damping avoids sinks from absorbing all of the mass and ensures primitivity of $G_{\mathrm{PR}}$ under mild conditions, guaranteeing uniqueness of the Perron vector \cend{23}.

\begin{figure}[htbp]
\centering
\begin{tikzpicture}[
  font=\small,
  vertex/.style={draw, circle, minimum size=0.95cm},
  arrow/.style={-{Stealth[length=2.5mm]}, thick},
  teleport/.style={arrow, dashed}
]
\node[vertex] (a) at (0,0) {$A$};
\node[vertex] (b) at (3.2,1.5) {$B$};
\node[vertex] (c) at (3.2,-1.5) {$C$};
\node[vertex] (d) at (6.4,0) {$D$};
% Graph transitions: A->B, A->C, B->D, C->D (same topology as the worked example; D is dangling)
\draw[arrow] (a) -- (b);
\draw[arrow] (a) -- (c);
\draw[arrow] (b) -- (d);
\draw[arrow] (c) -- (d);
% Uniform teleportation, drawn from D only: one dashed arrow to every node, including D itself.
% D->A runs through the gap between B and C, and D->B / D->C bend outward, so no dashed arrow crosses a solid edge.
\draw[teleport] (d) -- (a);
\draw[teleport] (d) to[bend right=32] (b);
\draw[teleport] (d) to[bend left=32] (c);
\draw[teleport] (d) to[loop right, min distance=12mm] (d);
\end{tikzpicture}
\caption{Random-surfer intuition for PageRank on the four-node graph of Section~\ref{sec:pagerank-worked-example}. Solid arrows are graph transitions, followed with total probability $d$. Dashed arrows are uniform teleportation, probability $(1-d)/n$ to every node including the current one; the dashed fan is drawn from $D$ only but leaves every node. $D$ has no outgoing edge, so under Equation~\eqref{eq:transition} its remaining mass is also spread uniformly, which is how damping keeps a sink from absorbing the chain.}
\label{fig:random-surfer}
\end{figure}

\dissertationSubsubheading[sub:power-iteration]{Power iteration and complexity}

In practice $\boldsymbol{\pi}$ is typically computed by power iteration:
\begin{equation}
\boldsymbol{\pi}^{(t+1)} = G_{\mathrm{PR}}\boldsymbol{\pi}^{(t)},
\label{eq:power-iteration}
\end{equation}
until $\|\boldsymbol{\pi}^{(t+1)}-\boldsymbol{\pi}^{(t)}\|_1 < \varepsilon$ or until a maximum iteration count is reached. Each iteration requires $O(|E|)$ arithmetic operations for sparse graphs \cend{49}. Methods that generate fewer edges, such as using latest-allowance snapshots rather than full transfer histories, therefore lower per-iteration cost and sometimes memory. This complexity fact motivates hypothesis H3.

Default parameter values used in this work follow the convention of setting damping $d=0.85$ \cend{15}, convergence tolerance $\varepsilon=10^{-8}$, and a maximum of $300$ iterations, as study-specific solver settings.\footnote{Tolerance and iteration caps are choices to facilitate reproducible run-times; they are not posited as part of the classical PageRank formulation.} The robustness section varies $d\in\{0.75,0.85,0.95\}$.

\dissertationSubheading[sec:reputation-systems]{Reputation systems lineage}

PageRank is one member of a family of spectral and propagation-based reputation mechanisms. Kleinberg's HITS algorithm scores pages both as hubs and authorities \cend{50}. EigenTrust takes the concept to peer-to-peer networks and normalizes local trust so that a malicious node cannot arbitrarily raise its own score \cend{16}. TrustRank begins with a seed set of trusted nodes and propagates their trust outward to suppress web spam \cend{17}. In each case, reputation is relational, which means that it is difficult to fake without paying for endorsements from nodes that have already earned reputation.

Sybil attacks continue to be the primary concern on-chain \cend{11}. It is inexpensive to create new wallets but more expensive to create wallets that receive allowances or transfers from different counterparties. The Sybil-adjusted stability proxies employed in this work (inbound counterparty ratio, tenure, active months) partially account for the diversity of interactions but do not prevent a malicious actor from manipulating them. Chapter~\ref{ch:discussion} provides a cost model for Sybils in allowance graphs.

While machine-learning credit models using DeFi features \cend{14} can provide better predictive accuracy on labeled defaults, they sacrifice interpretability. EndorseRank and AWP focus on interpretable, audit-friendly propagation rules \cend{2}.

\dissertationSubheading[sec:allowance-endorsement]{ERC-20 allowances as endorsement}

The ERC-20 \texttt{approve} function and EIP-2612 \texttt{permit} allow users to explicitly authorize spending limits on their behalf \cend{18,19}. EndorseRank uses the latest allowance from an owner to a spender to represent the edge strength on a directed endorsement graph and applies PageRank to that snapshot. Since newer approvals supersede older ones, this does not require temporal decay.

Specifically, for each token, for each owner--spender pair, let $a_{o\to s}$ denote the decimal-adjusted latest allowance. Across all tokens, we aggregate allowances by summing the individual allowances to obtain a single edge weight $w_{o\to s}$. The resulting graph $G_{\text{approve}}=(V,E_{\text{approve}},w)$ is usually much sparser than the transfer graph over the same wallets and window. PageRank on $G_{\text{approve}}$ gives EndorseRank scores $\mathbf{s}_{\mathrm{ER}}$.

This endorsement interpretation is deliberate but imprecise. Users approve routers because it is convenient to do so, not necessarily because they want to make a social statement about the trustworthiness of the router. Contract addresses and EOAs coexist. But allowances are a more explicit signal of approval than receiving transfers, which could be market making, wash trading \cend{29}, airdrop farming \cend{32}, or routing without any endorsement intent.

\dissertationSubheading[sub:gf-pr]{What this research does not claim}

This research does not claim that EndorseRank predicts trading profit, loss avoidance or liquidation. The research proposal intended to validate both methods against proxies for success at trading on GMX V2, and to bound that validation with an outcome-native reference (PageRank on a star graph constructed from realized profit and loss). This reference is circular by definition, as it is derived from the same PnL stream that is used to define the trading-success proxies, so its almost perfect agreement with the realized-gain proxy is simply the correlation of a variable with a monotone transformation of itself, rather than information about reputation. By the same reasoning, we cannot infer that a social method has failed when it weakly agrees with those proxies. The reference, its accompanying supplementary research question and claim, and the trading-success, inverse-risk and liquidation families are consequently removed from the main comparison; their matrices can be found in Appendix~\ref{ch:appendix-archive}. The only other external evaluation is a temporal holdout of future new owner--spender approvals (Chapters~\ref{ch:methodology} and~\ref{ch:results}), which is out-of-window but the same construct as the score.

\dissertationSubheading[sec:rank-correlation-theory]{Rank correlation theory}

Reputation scores and validation proxies are continuous variables without a natural classification boundary. Building on the AWP validation design\cend{3} and classical psychometrics \cend{22}, this research assesses monotonic and pairwise rank agreement, and estimates the sampling uncertainty of each reported coefficient with bootstrap confidence intervals \cend{34}.

Spearman's rank correlation \cend{20} is Pearson correlation on ranks. For paired ranks $(R_i,S_i)$ of $n$ items,
\begin{equation}
\rho = 1 - \frac{6\sum_{i=1}^n (R_i-S_i)^2}{n(n^2-1)}
\label{eq:spearman}
\end{equation}
in the absence of ties (with tie-aware variants used in software). Spearman's $\rho$ captures monotonic association.

Kendall's $\tau$ \cend{21} counts concordant and discordant pairs. For all pairs $(i,j)$ with $i<j$,
\begin{equation}
\tau = \frac{2(C - D)}{n(n-1)},
\label{eq:kendall}
\end{equation}
where $C$ ($D$) is the number of concordant (discordant) pairs, with tie corrections in the $\tau$-b variant. Kendall's $\tau$ is usually easier to interpret for agreement about pairwise orderings and is the main family-level summary in Chapter~\ref{ch:results}. The values are compared relative to other methods on the same cohort, not to an uncited absolute value.

\dissertationSubheading[sec:theoretical-framework]{Theoretical framework}

This study uses hyperlink--citation propagation, construct (domain) alignment, and computational complexity as the bases for comparing EndorseRank and AWP across reputation structure (H1), social-method alignment (H2), and computational efficiency (H3).

\begin{itemize}
\item Hyperlink--citation semantics: The original PageRank algorithm \cend{15} treats inbound hyperlinks as citations. Do and Do\cend{4} apply this logic to Ethereum transfers; AWP\cend{3} uses time-decayed transfer edge weights. Figure~\ref{fig:logistic-decay-methods} (Chapter~\ref{ch:methodology}) plots the logistic decay
\begin{equation}
\sigma(\Delta t)=\frac{1}{1+\exp\bigl(k(\Delta t-t_0)\bigr)}
\label{eq:logistic}
\end{equation}
with $k=0.01$ and $t_0=180$ days as used for AWP in this study. EndorseRank extends citation-based propagation to ERC-20 approve/permit edges because allowances are explicit, on-chain delegations of spending authority \cend{18,19}. Unlike transfers, which may reflect trading or routing without trust intent, an allowance records explicit authorization that remains in force until revoked or replaced. Graph-based risk models \cend{12} further show that network structure carries information beyond raw activity counts, which supports a directed endorsement graph in addition to transfer volume.

\item Domain alignment framework: In measurement theory, construct alignment assesses how well an indicator tracks an intended construct \cend{22}; Messick\cend{33} extends this to a unified judgment about what a score means and what its use implies. For rank-based reputation scoring, Spearman's $\rho$ and Kendall's $\tau$ evaluate monotonic and pairwise ordering agreement without assuming a classification threshold, as in the AWP validation protocol\cend{3}. Contemporaneous transfer and allowance statistics check graph coherence; a temporal holdout of future new approvals is the out-of-window check.

\item Computational complexity: PageRank's per-iteration cost is $O(|E|)$ \cend{15}. Allowance graphs are sparser than transfer graphs (fewer edges $|E|$), so per-iteration cost and memory are lower and, in practice, fewer iterations are often needed; this is the basis for H3. Because the solver is identical for both methods, any runtime difference is an expected consequence of edge sparsity, not an algorithmic contribution.
\end{itemize}

Figure~\ref{fig:conceptual-model} presents the conceptual model mapping the primary EndorseRank--AWP comparison to H1--H3 and RQ1--RQ4; RQ5 sits on top of it as the question of whether the two constructs combine.

\begin{figure}[htbp]
\centering
% !TEX root = ../main.tex
\begin{tikzpicture}[
  font=\small,
  box/.style={draw, rounded corners, align=center, minimum width=10.5cm, minimum height=1.0cm},
  smallbox/.style={draw, rounded corners, align=center, minimum width=3.4cm, minimum height=1.05cm},
  refbox/.style={draw, dashed, rounded corners, align=center, minimum width=3.4cm, minimum height=0.95cm},
  arrow/.style={-{Stealth[length=2.5mm]}, thick}
]
\node[box] (center) at (0,3.2) {Primary comparison: EndorseRank vs AWP};

\node[smallbox] (rep) at (-5,1.2) {Reputation\\Structure};
\node[smallbox] (risk) at (0,1.2) {Construct\\Checks};
\node[smallbox] (eff) at (5,1.2) {Computational\\Efficiency};

\node[smallbox] (h1) at (-5,-1.0) {H1 / RQ1\\EndorseRank--AWP\\divergence};
\node[smallbox] (h2) at (0,-1.0) {H2 / RQ2\\Allowance vs\\transfer};
\node[smallbox] (h3) at (5,-1.0) {H3 / RQ3\\EndorseRank\\efficiency};

\node[refbox] (hold) at (0,-3.0) {RQ4: temporal holdout\\future new approvals};

\draw[arrow] (center.south -| rep.north) -- (rep.north);
\draw[arrow] (center.south) -- (risk.north);
\draw[arrow] (center.south -| eff.north) -- (eff.north);

\draw[arrow] (rep.south) -- (h1.north);
\draw[arrow] (risk.south) -- (h2.north);
\draw[arrow] (eff.south) -- (h3.north);

\draw[arrow, dashed] (h2.south) -- (hold.north);
\end{tikzpicture}

\caption{Conceptual model: primary EndorseRank vs.\ AWP comparison mapped to reputation structure (H1/RQ1), construct checks (H2/RQ2), computational efficiency (H3/RQ3), and the temporal holdout of future new approvals (RQ4).}
\label{fig:conceptual-model}
\end{figure}

H1 and H3 concern EndorseRank and AWP only; H2/RQ2 compares contemporaneous allowance and transfer checks; RQ4 is a time-split test of future new owner--spender approvals on the spender cohort, not a ranking of trading outcomes.

\dissertationSubheading[sub:terminology]{Terminology}

The definitions in Chapter~\ref{ch:introduction} establish formal constructs. The following terms are used when interpreting results:

\begin{itemize}
\item Primary methods: EndorseRank (allowance graph) and AWP (transfer graph). Both are social PageRank methods on wallet-to-wallet edges; intended as deployable reputation signals.
\item Primary graph construct: The graph a method is \emph{designed} to rank (e.g., $G_{\text{approve}}$ for EndorseRank). Distinct from checks applied after scoring.
\item Contemporaneous check: Local transfer or allowance statistics computed in the same window as the score \cend{22}.
\item Temporal holdout: Scores frozen at t1; labels from t2 new owner--spender approvals on inbound spenders.
\item Matched wallet cohort: The primary same-window sample of $n=5{,}521$ wallets with non-zero endorsement- and transfer-graph connectivity.
\item Spender holdout cohort: $n=1{,}335$ t1 inbound allowance recipients used for RQ4.
\item Expanded wallet pool: A larger address set ($N=99{,}080$) used for the node-count sensitivity of runtime beyond the matched cohort.
\end{itemize}

\dissertationSubheading[sec:domain-alignment-metrics]{Domain alignment metrics}

Following the AWP validation protocol\cend{3} and Cronbach and Meehl\cend{22}, construct alignment for rank-based scores uses Spearman's $\rho$ and Kendall's $\tau$ between reputation scores and check rankings. This research reports contemporaneous transfer and allowance checks and a temporal holdout of future new approvals; see Chapter~\ref{ch:methodology} for operational definitions.

The literature reviewed here supports three design choices: (i) PageRank as an interpretable propagation backbone; (ii) allowances as endorsement edges distinct from transfers; and (iii) rank correlation plus a time split as the check design. These choices motivate the remaining sections of this chapter and, later, the operational pipeline in Chapter~\ref{ch:methodology}.

\dissertationSubheading[sec:pagerank-worked-example]{Worked example: PageRank on a four-node endorsement graph}

Consider four wallets $\{A,B,C,D\}$ with latest-allowance edges
\[
A\to B\ (2),\quad A\to C\ (1),\quad B\to D\ (3),\quad C\to D\ (1),
\]
and damping $d=0.85$. Out-strengths are $d_A^{\mathrm{out}}=3$, $d_B^{\mathrm{out}}=3$, $d_C^{\mathrm{out}}=1$, $d_D^{\mathrm{out}}=0$ (dangling). After one power-iteration step from the uniform vector $\boldsymbol{\pi}^{(0)}=\mathbf{1}/4$, mass flows preferentially toward $D$ because both $B$ and $C$ endorse $D$, while $A$ splits mass between $B$ and $C$. Subsequent iterations mix in teleportation mass $(1-d)/4=0.0375$ at every node. The stationary ranking places $D$ highest, then $B$ and $C$, then $A$. The example shows that \emph{receiving} endorsements from nodes that themselves receive endorsements raises score, which is exactly the citation logic of PageRank \cend{15,23}.

If the same wallets are connected by many small transfers, rather than a few large ones, then the transition matrix will be different, and we might see a different ranking even though the nodes are the same. This is essentially the microcosm of hypothesis H1: EndorseRank and AWP need not agree.

\dissertationSubheading[sec:defi-risk-taxonomy]{DeFi risk taxonomy and where reputation fits}

Werner et al.\cend{8} categorize DeFi risk by protocol layer, while Zhou et al.\cend{27} organize it by the types of attacks observed in practice. A bug in the contract code or a misconfigured upgrade key is a smart-contract risk \cend{51}. An outdated price or a feed that someone has controlled is an oracle risk \cend{39}. Market volatility and liquidation cascades, including stress in stablecoin and funding markets \cend{53}, are market risks \cend{52}; the market crash in March 2020 showed how these two layers can interact \cend{9}. Token-weighted voting can be modeled, and that is a governance risk \cend{54}. In this paper, we study a more specific layer, the \emph{counterparty} or \emph{behavioral} risk: whether a wallet's historical behavior suggests that it manages leveraged exposure responsibly.

On-chain reputation scores only capture part of the behavioral risk layer, and they do not replace collateral, auditing, or oracle design. Rather, their purpose is to condense the relational structure of a system (who authorizes whom, who pays whom) into an ordinal measure that is easy to audit and reproduce \cend{2}. EndorseRank focuses on the authorization aspect of behavioral risk; AWP focuses on the payment-flow aspect. Neither aims to predict exploits or governance attacks.

Finally, fake activity makes measuring behavioral risk harder. Wash trading increases the trading volume on decentralized exchanges \cend{28} and, at a larger scale, on unregulated centralized exchanges \cend{29}; front-running and similar forms of extractable value produce transfers that are driven by ordering strategies, not relationships \cend{55}, and Qin et al.\cend{30} estimate how much of this value is captured on Ethereum. Uniswap scam tokens \cend{31} and cross-exchange arbitrage trades \cend{56} similarly increase the number of transactions but are not sustained relationships. The transfer graph is particularly sensitive to fake volume; the allowance graph is sensitive to approval farming. Time-based splits and Sybil detection are partial solutions, not full ones.

\dissertationSubheading[sec:related-measurement]{Related measurement traditions}

Reputation systems have been explored beyond the blockchain context, in particular in e-commerce \cend{57}, peer-to-peer networks \cend{16}, and web spam \cend{17}. The main insights that apply here are: (i) reputation is a property of relations; (ii) easy identity creation allows Sybil attacks \cend{11}; (iii) seeding and normalization reduce the impact of inflation; and (iv) evaluation must go beyond mere internal coherence, using external benchmarks \cend{22}. EndorseRank follows points (i)--(iii) through PageRank damping and endorsement costs, and point (iv) through contemporaneous construct checks, measured using Spearman's $\rho$ \cend{20} and Kendall's $\tau$ \cend{21}, as well as a temporal holdout.

Personalized and topic-sensitive variants of PageRank \cend{58,59} skew the teleportation vector toward a specific group of nodes; it is this mechanism that allows TrustRank to transform a tiny initial set of nodes into a global ranking score \cend{17}. Two separate bodies of work generalize PageRank to graphs containing multiple types of edges. In Multiplex PageRank, the centrality value of a node on one layer affects the random walk on another layer \cend{64}; in the multilayer framework proposed by De Domenico et al., a single random walker moves both within and between layers of a single interconnected network \cend{65}. Both approaches model layers as alternative views of a single set of nodes, which is what we have here: allowance edges and transfer edges represent two different relationships between the same set of addresses. Both EndorseRank and AWP use a uniform teleportation vector and only a single layer to allow direct comparison with our baseline \cend{3}. The combined operator described in Chapter~\ref{ch:methodology} has the simplest form of multilayer PageRank possible, namely a weighted sum of the two transition matrices per node using a single weight, and the seeded ablation experiments make use of the personalized teleportation idea by using the allowance score as the teleportation vector.

\dissertationSubheading[sec:account-model]{Accounts, gas, and why logs are the measurement surface}

Accounts in Ethereum are of two kinds, externally owned accounts (EOAs) and contract accounts. EOAs are owned by a private key; contract accounts are owned by the code they store \cend{60}. All transactions that alter state consume gas and generate logs when contract accounts write to storage. If reputation systems intend to be auditable, then they must prefer logs over off-chain databases: logs are a permanent part of the blockchain history and can be independently recomputed by any full node or index provider \cend{61}.

Gas prices also influence the structure of the resulting graphs. On L2 networks, cheaper gas leads to higher transfer volume and more approval churn, densifying $G_{\text{transfer}}$ more than $G_{\text{approve}}$ because transfer activity is a regular occurrence while approval activity happens once in order to permit future transfer activity. It is for this reason that we expect the graph produced by EndorseRank to be considerably less dense than the transfer graph, an observation we explore further in Chapter~\ref{ch:results}.

\dissertationSubheading[sec:approval-security]{Approvals as security-relevant events}

An unlimited approval (\texttt{type(uint256).max}) is a well-known user-experience issue: if the recipient account is compromised, all of the owner's tokens can be stolen until the owner resets the approval. Allowances are considered capabilities rather than mere metadata when analyzing honeypot smart contracts \cend{62} and when analyzing phishing accounts in Ethereum \cend{24}. EndorseRank considers a token approval as endorsing a particular spender's capability to transfer those tokens, which is why an allowance is not semantically equivalent to a transfer, even though both involve the movement of tokens. An account could receive many transfers but never be granted permission to transfer funds on behalf of others.

From a data perspective, unlimited approvals introduce extreme edge weights. The decimal adjustment, and the optional winsorization, are there to dampen outliers; this work combines the decimal-adjusted latest allowances, and uses PageRank's normalization by out-strength such that a single very large allowance will not always win, only if the account has few outgoing endorsements.

\dissertationSubheading[sec:awp-in-depth]{AWP in depth: what the transfer-graph baseline measures}

Do and Do\cend{4} use PageRank and HodgeRank as signals of creditworthiness by applying them to Ethereum transfer graphs. Adaptive Weighted PageRank (AWP), proposed by Do, Do, and Nguyen\cend{3}, uses token-transfer edges whose weight is value and recency via logistic decay, and validates ranking accuracy using inbound degree and inbound value (as proxies for ``the amount of economic attention a wallet attracts''). The decay function (Equation~\eqref{eq:logistic}) and the in-degree/in-value proxies used in Chapters~\ref{ch:methodology} and~\ref{ch:results} both originate from the AWP paper. This approach is appealing since it is entirely on-chain, ranking based, and not dependent on any kind of off-chain identities.

EndorseRank continues the template of PageRank plus a rank-correlation metric, but changes the edge semantics. In Chapter~\ref{ch:results} we investigate whether allowance edges produce a construct-distinct, allowance-aligned, and less expensive score, and whether the t1 scores can predict new approvals at t2.

\dissertationSubheading[sec:method-comparison-table]{Comparison of graph reputation methods}

Table~\ref{tab:method-comparison-lit} places EndorseRank alongside other approaches covered in this chapter, categorized by edge semantics, temporal model and validation design.

\begin{table}[htbp]
\centering
\caption{Comparison of selected graph reputation methods (conceptual).}
\label{tab:method-comparison-lit}
\fitwidth{%
\begin{tabular}{p{0.18\linewidth}p{0.24\linewidth}p{0.24\linewidth}p{0.28\linewidth}}
\toprule
Method & Edge semantics & Temporal model & Typical validation \\
\midrule
PageRank (web) \cend{15} & Hyperlinks & None / static snapshot & Retrieval quality \\
EigenTrust \cend{16} & Local trust ratings & Iterative aggregation & P2P file authenticity \\
TrustRank \cend{17} & Links from seed goods & Static / seeded & Spam demotion \\
PageRank/HodgeRank on Ethereum \cend{4} & Token transfers & Static window & Not reproduced in this thesis \\
AWP \cend{3} & Token transfers & Logistic decay & In-degree / in-value \\
RiskProp \cend{12} & Risk-bearing links & Propagation of risk & Account risk labels \\
Multiplex / multilayer PageRank \cend{64,65} & Several relations over one node set & Static layers & Versatility of top-ranked nodes \\
EndorseRank (this study) & Latest allowances & Snapshot (override) & Allowance checks; t2 new approvals \\
C-PR / S-PR (this study) & Allowances mixed with, or seeding, transfers & Snapshot and logistic decay & Both holdout labels; pre-registered rule \\
\bottomrule
\end{tabular}
}
\end{table}

Table~\ref{tab:method-comparison-lit} indicates that EndorseRank's innovation lies in selecting the latest-allowance endorsement edge as the main social construct for DeFi wallets rather than the ranking algorithm itself. Our contribution in this regard is altering the operator: specifically, we couple the operators, and perform a seeded ablation to see how much this coupling contributes, by merging the two relations into one walk; the test of this coupling is out of window, against both constructs, and under a pre-defined rule.

\dissertationSubheading[sec:literature-summary-table]{Summary of the reviewed literature}

Table~\ref{tab:literature-summary} provides a consolidated list of the sources on which this study is based. For each source it notes the methodology, the strengths that this thesis leverages, the limitations that this thesis attempts to address, and the relation to this work. This table constitutes the evidentiary basis for the research gap stated in Section~\ref{sec:research-gap}.

{\footnotesize
\setlength{\tabcolsep}{4pt}
\begin{longtable}{p{0.15\linewidth}p{0.20\linewidth}p{0.19\linewidth}p{0.19\linewidth}p{0.18\linewidth}}
\caption{Summary of the reviewed literature: methodology, strengths, limitations, and relation to this study.}
\label{tab:literature-summary}\\
\toprule
Source (year) & Methodology & Strengths & Limitations & Relation to this study \\
\midrule
\endfirsthead
\multicolumn{5}{l}{\small Table~\thetable\ (continued)}\\
\toprule
Source (year) & Methodology & Strengths & Limitations & Relation to this study \\
\midrule
\endhead
\midrule
\multicolumn{5}{r}{\small continued on next page}\\
\endfoot
\bottomrule
\endlastfoot
Page et al.\ (1998) \cend{15}; Langville and Meyer (2006) \cend{23} & Eigenvector ranking of a hyperlink graph; power iteration & Interpretable propagation; $O(|E|)$ per iteration; unique Perron vector under damping & Edge semantics fixed to hyperlinks; no notion of time & Solver and complexity argument reused unchanged for both compared methods \\
Kamvar et al.\ (2003) \cend{16}; Gy\"{o}ngyi et al.\ (2004) \cend{17} & Trust aggregation in P2P networks; seeded propagation against spam & Normalisation and seed sets limit inflation; explicit adversary model & Requires explicit ratings or curated seeds; not on-chain & Motivates the Sybil-cost discussion and the uniform-teleportation baseline \\
Do and Do (2023) \cend{4} & PageRank and HodgeRank on Ethereum transaction graphs as social credit & Framing of transfer connectivity as credit-relevant; fully public data & Transfers conflate payment and trust; no allowance semantics & Prior work on the transfer graph; distinguished from AWP in this thesis \\
Do, Do, and Nguyen (2023), AWP \cend{3} & Value- and recency-weighted transfer PageRank with logistic decay; rank correlation against in-degree/in-value & Rank-based validation without default labels; recency handled explicitly & Proxies are computed from the same transfer data that builds the graph (mechanical alignment); dense graph & Baseline method, decay formula, and proxy protocol reproduced exactly on the same cohort \\
Lin et al.\ (2025), RiskProp \cend{12} & De-anonymisation score plus network propagation of account risk & Shows structure carries information beyond activity counts; labelled evaluation & Needs risk labels; targets fraud, not trust or authorisation & Motivates refined edge semantics; RiskProp-style propagation kept as an archived extended baseline \\
Wu et al.\ (2022) \cend{24} & Network embedding of Ethereum transaction graphs for phishing detection & Graph features outperform per-account features; large labelled set & Supervised; labels are phishing, not creditworthiness & Supports the premise that graph position is informative; not a competitor method \\
Hassija et al.\ (2020) \cend{13}; Kotb (2023) \cend{14} & Prospect-theory and machine-learning credit scoring on blockchain data & Directly target lending decisions & Require labels or off-chain features; limited interpretability & Contrast class: EndorseRank stays within auditable graph propagation \\
Bellini et al.\ (2020) \cend{25}; Almasoud et al.\ (2020) \cend{26} & Surveys of blockchain-based trust and reputation systems & Map the design space; identify explicit-rating and volume-based signals as dominant & Do not evaluate authorisation edges; no L2 measurement & Establish that latest-allowance edges are absent from the surveyed designs \\
Victor (2020) \cend{40}; Victor and Weintraud (2021) \cend{28}; Cong et al.\ (2023) \cend{29} & Address clustering heuristics; wash-trading detection on DEXs and exchanges & Quantify how much on-chain activity is artificial or multi-address & Do not propose a reputation score & Explain why transfer-based signals are exposed to inflation; motivate the allowance construct \\
Cronbach and Meehl (1955) \cend{22}; Messick (1995) \cend{33} & Construct validity in psychometrics & Separate what a score is designed to measure from external criteria & Not blockchain-specific & Frame the contemporaneous checks as construct checks and the temporal holdout as external evidence \\
Efron (1979) \cend{34}; DiCiccio and Efron (1996) \cend{35} & Bootstrap resampling and bootstrap confidence intervals & Distribution-free uncertainty for rank statistics; paired contrasts & Assumes exchangeable units & Provides the 95\% intervals and $\Delta\tau$ contrasts reported in Chapter~\ref{ch:results} \\
\end{longtable}
}

\dissertationSubheading[sec:research-gap]{Research gap}

There are five gaps derived from Table~\ref{tab:literature-summary} and Section~\ref{sec:related-measurement}. First, all of the on-chain reputation methods discussed (PageRank/HodgeRank \cend{4}, AWP \cend{3}, RiskProp \cend{12}, and those found in the surveys \cend{25}) build their graph from transfers or explicit ratings; none makes use of the ERC-20 allowance, the only on-chain event that represents a conscious delegation of the power to spend. Whether the construction of a latest-allowance endorsement graph leads to a distinct and meaningful ordering is unknown (O1 and RQ1). Second, the transfer-graph baseline's validation process compares scores against in-degree and in-value calculated using the same transfer data that was used to build the graph, meaning that some agreement is inevitable; there is no report of such a process being applied to an allowance construct and its transfer equivalent using confidence intervals and paired difference tests on the same set of addresses (O2 and RQ2). Third, runtimes for blockchain PageRank are often given without information about the number of edges, amount of memory used, and number of iterations made using a fixed protocol, making it difficult to attribute any claims regarding efficiency to sparsity rather than implementation (O3 and RQ3). Fourth, there is a methodological gap related to time: graph statistics gathered during the same window are often given as validation despite the fact that such scores cannot be shown to predict future behavior; instead, the out-of-window check used here involves a temporal holdout with a future-tense label for each construct, and the raw degree at the time of the freeze being used as the baseline (O4 and RQ4). Fifth, although multilayer ranking has been developed for social, transport and citation networks \cend{64,65}, it has not been used on two on-chain relations over one set of addresses; no study provides evidence of whether using a combination of an authorization layer and a payment layer performs better than either layer individually when the rules are defined before seeing the outcome (O5 and RQ5).

\dissertationSubheading[sec:open-problems]{Open problems motivating this research}

Several open problems in on-chain reputation remain unresolved in the literature and motivate the present study:

\begin{enumerate}
\item Edge semantics: Which on-chain relations best encode trust or authorization under pseudonymity \cend{2}?
\item Validation without defaults: How can methods be compared when loan-default labels are scarce \cend{4}?
\item Computational cost: Can reputation scores refresh at interactive latencies as graphs grow \cend{23}?
\item Time-split checks: How can evaluation avoid mistaking same-window graph statistics for out-of-window evidence \cend{22}?
\item Adversarial robustness: What is the cost of forging reputation under cheap identities \cend{11}?
\end{enumerate}

EndorseRank addresses (1)--(4) with allowance edges, contemporaneous checks, runtime benchmarks, and a temporal holdout, and analyzes (5) at a model level in Chapter~\ref{ch:discussion}. Full empirical adversarial evaluation remains future work.

Chapter~\ref{ch:methodology} turns the gap statement into an operational design: it fixes the data window and cohort, specifies how both graphs are constructed and solved with one shared PageRank implementation, defines the proxy families and the temporal holdout, and sets the timing and bootstrap protocols whose outputs Chapter~\ref{ch:results} reports.

